\documentclass[10pt, conference]{IEEEtran}
\usepackage{cite}
\usepackage{amsmath,amssymb,amsfonts}
\usepackage{algorithmic}
\usepackage{graphicx}
\usepackage{textcomp}
\usepackage{xcolor}
\usepackage{booktabs}
\usepackage{url}

\begin{document}

\title{Aeon: Training-Free Temporal and Semantic Enhancements for Hyper-Realistic Visualizations of Ancient Egypt via Video Diffusion}

\author{\IEEEauthorblockN{Negma Abdelrahman, Ramez Ezzat, Mariam Mohamed, Marc Hany, and Rokia Islam}
\IEEEauthorblockA{Department of Artificial Intelligence Science\\
Alamein International University\\
El Alamein, Egypt}
}

\maketitle

\begin{abstract}
Recent advancements in video generation have been driven by large-scale open-weight Diffusion Transformers (DiTs). However, these models frequently suffer from temporal instability, geometric distortion, and weak semantic adherence, especially in historically rich themes like Ancient Egypt. In this paper, we propose \textbf{Aeon}, a zero-cost, training-free methodology to enhance the temporal consistency and semantic alignment of open-weight video diffusion models. Evaluating our approach on the 10-billion parameter Mochi-1 architecture, we isolate artifacts through three inference-time modifications: (1) shifting the Flow Match scheduler distribution, (2) injecting structured asymmetric negative embeddings, and (3) amplifying classifier-free guidance. Our method improves structural similarity (SSIM) by 1.3\% and perceptual quality (LPIPS) by 21\% without incurring additional computational overhead, enabling hyper-realistic visualizations of ancient desert landscapes and monuments.
\end{abstract}

\begin{IEEEkeywords}
Video Generation, Diffusion Models, Flow Matching, Ancient Egypt, Zero-Shot Learning
\end{IEEEkeywords}

\section{Introduction}
Text-to-video (T2V) generation has seen rapid progress with the scaling of Diffusion Transformers (DiTs). Open-weight models like Mochi-1-preview represent the state-of-the-art in accessible video generation. However, achieving hyper-realism in historically demanding themes, such as "Ancient Egypt," remains a significant challenge. The complexity of desert lighting, the rigidity of monumental architecture like pyramids, and the fluid motion of sand require exceptional temporal and semantic stability.

In this work, we introduce \textbf{Aeon}, a training-free enhancement pipeline designed to bring Ancient Egypt to life with unprecedented realism. We identify and mitigate artifacts through native inference-time modifications, requiring no additional training or fine-tuning. Our contributions include:
\begin{itemize}
    \item A mathematical analysis of the diffusion reverse process and flow matching velocity.
    \item Identification of "Temporal Morphing" as a primary weakness in rigid geometric structures.
    \item A zero-cost methodology to enhance structural similarity by 1.3\% and perceptual quality by 21\%.
\end{itemize}

\section{Background and Mathematical Foundation}
\subsection{Denoising Diffusion Probabilistic Models (DDPM)}
The diffusion process consists of a forward Markov chain $q(\mathbf{x}_{1:T} | \mathbf{x}_0)$ that gradually adds Gaussian noise to data $\mathbf{x}_0$ according to a variance schedule $\beta_1, \dots, \beta_T$:
\begin{equation}
q(\mathbf{x}_t | \mathbf{x}_{t-1}) = \mathcal{N}(\mathbf{x}_t; \sqrt{1-\beta_t}\mathbf{x}_{t-1}, \beta_t \mathbf{I})
\end{equation}
The reverse process $p_\theta(\mathbf{x}_{t-1} | \mathbf{x}_t)$ aims to reconstruct the original data by learning to denoise:
\begin{equation}
p_\theta(\mathbf{x}_{t-1} | \mathbf{x}_t) = \mathcal{N}(\mathbf{x}_{t-1}; \mathbf{\mu}_\theta(\mathbf{x}_t, t), \mathbf{\Sigma}_\theta(\mathbf{x}_t, t))
\end{equation}

\subsection{Score Function and Training Objective}
The score function $\nabla_{\mathbf{x}} \log p(\mathbf{x})$ points towards higher density regions. In diffusion models, learning to denoise is equivalent to learning the score function. The training objective is derived from the Evidence Lower Bound (ELBO), which simplifies to a Mean Squared Error (MSE) loss in the noise space:
\begin{equation}
L_{simple} = \mathbb{E}_{t, \mathbf{x}_0, \epsilon} [ \| \epsilon - \epsilon_\theta(\mathbf{x}_t, t) \|^2 ]
\end{equation}

\subsection{Noise Schedules and Sampling}
The noise schedule (linear vs. cosine) dictates the rate of information loss. Cosine schedules provide better stability by avoiding abrupt information collapse. During inference, sampling methods like DDPM require $T$ steps, whereas Denoising Diffusion Implicit Models (DDIM) enable faster deterministic sampling by bypassing the Markovian assumption.

\section{Model Architecture}
\begin{figure}[htbp]
\centering
\includegraphics[width=\linewidth]{architecture_diagram.png}
\caption{System architecture of Aeon. The pipeline utilizes a VAE for latent compression, a T5-XXL text encoder for conditioning, and a stack of Asymmetric Diffusion Transformer (AsymmDiT) blocks with spatial-temporal attention.}
\label{fig:architecture}
\end{figure}

\subsection{Asymmetric Diffusion Transformer (AsymmDiT)}
Our system utilizes Mochi-1's Asymmetric Diffusion Transformer (AsymmDiT). Unlike standard 3D U-Nets, AsymmDiT treats video data as a sequence of latent patches, incorporating spatial-temporal self-attention blocks to capture long-range dependencies across both dimensions.

\subsection{Latent Space and Conditioning}
Diffusion occurs in a compressed latent space provided by a Variational Autoencoder (VAE). Conditioning is achieved through cross-attention with text embeddings from a T5-XXL encoder. Temporal consistency is enforced via 3D positional encodings and temporal attention layers that model frame-to-frame transitions.

\section{Loss Function Analysis}
Mochi-1 utilizes a Flow Matching objective, predicting the velocity $\mathbf{v}_t$ (the derivative of the latent path):
\begin{equation}
L_{flow} = \mathbb{E}_{t, \mathbf{x}_0, \epsilon} [ \| \mathbf{v}_t - \mathbf{v}_\theta(\mathbf{x}_t, t) \|^2 ]
\end{equation}
Velocity prediction ($\mathbf{v}$) is superior to noise prediction ($\epsilon$) in high-resolution video tasks as it provides a stable training signal across the entire ODE trajectory. We justify the use of $v$-prediction for its ability to preserve structural geometry during complex motion.

\section{Weakness Analysis}
In the context of "Ancient Egypt's Realism," we identified two primary weaknesses in the baseline model:
\begin{enumerate}
    \item \textbf{Geometric Instability:} Rigid structures like pyramids exhibit "melting" textures during pans, caused by the scheduler's failure to resolve structural details early in the sampling process.
    \item \textbf{Semantic Drift:} Prompts involving specific Egyptian lighting (e.g., "sunset over Luxor") lose consistency, indicating a decay in guidance adherence in later diffusion steps.
\end{enumerate}

\section{Proposed Enhancements}

\begin{table*}[t]
\centering
\caption{Ablation Study Results (Mean $\pm$ Std). Best results in bold.}
\label{tab:ablation_results}
\begin{tabular}{lccccc}
\toprule
Config & CLIP-SIM & SSIM & LPIPS & Time (s) & FVD vs Baseline \\
\midrule
Baseline & 32.09 $\pm$ 2.10 & 0.9393 $\pm$ 0.0715 & 0.0522 $\pm$ 0.0514 & 492.4 $\pm$ 31.8 & N/A (reference) \\
A: Scheduler Only & 31.98 $\pm$ 2.00 & 0.9487 $\pm$ 0.0462 & 0.0462 $\pm$ 0.0421 & 488.5 $\pm$ 31.7 & 90.87 $\pm$ 71.82 \\
B: Neg Prompt Only & 32.67 $\pm$ 2.47 & 0.9509 $\pm$ 0.0498 & 0.0439 $\pm$ 0.0412 & 481.8 $\pm$ 29.6 & 145.66 $\pm$ 85.90 \\
C: Guidance Only & \textbf{33.13 $\pm$ 2.37} & 0.9404 $\pm$ 0.0679 & 0.0575 $\pm$ 0.0699 & \textbf{476.0 $\pm$ 30.5} & 150.53 $\pm$ 94.31 \\
D: All Combined & 32.65 $\pm$ 2.70 & \textbf{0.9511 $\pm$ 0.0516} & \textbf{0.0410 $\pm$ 0.0413} & 476.3 $\pm$ 29.1 & 136.36 $\pm$ 76.92 \\
\bottomrule
\end{tabular}
\end{table*}
Our enhancements are applied directly within the reverse diffusion loop of the Mochi-1 pipeline.

\subsection{Scheduler Optimization (Shift Adjustment)}
Mochi-1 utilizes a Flow Match Euler Discrete Scheduler. The \texttt{shift} parameter dictates the distribution of timesteps sampled during inference. By default, Mochi utilizes a shift of 1.0. We increase this parameter to 1.15 to prioritize the resolution of structural geometry and high-frequency details. 

\subsection{Structured Negative Prompt Injection}
Early video models frequently struggle with object permanence. We inject a structured negative prompt: \textit{"blurry, morphed, morphing, jittery, low resolution, deformed, unnatural motion, artifacts, bad anatomy"}. The diffusion trajectory is actively pushed away from these artifact-associated latent regions.

\subsection{Guidance Amplification}
To complement the negative prompt, we increase the CFG scale from 4.5 to 6.0. This stronger guidance enforces stricter adherence to the positive Egyptian semantic prompt while penalizing deviations toward negative embeddings.

\section{Experimental Setup}
\subsection{Evaluation Prompts}
To ensure robustness, we evaluate our pipeline on five diverse prompts centered on Ancient Egyptian themes:
\begin{enumerate}
    \item \textit{"A camel moving in the desert with pyramids in the background, dramatic sunset"}
    \item \textit{"Ancient Egyptian priests walking through a temple corridor, flickering torches"}
    \item \textit{"The Nile river at dawn with ancient sailboats and lush greenery, cinematic lighting"}
    \item \textit{"A pharaoh standing on a balcony overlooking a bustling ancient city, high detail"}
    \item \textit{"A close-up of an ornate golden sarcophagus in a dimly lit tomb, dust motes in the air"}
\end{enumerate}

\subsection{Metrics}
We evaluate our pipeline using four metrics that capture the trade-off between semantic adherence and temporal stability:
\begin{itemize}
    \item \textbf{CLIP-SIM}: Measures the cosine similarity between the CLIP embeddings of the generated frames and the text prompt. \textit{Justification}: It is the gold standard for quantifying how well the model adheres to the Egyptian thematic conditioning.
    \item \textbf{SSIM (Structural Similarity Index)}: A pixel-level metric that measures the structural consistency between consecutive frames. \textit{Justification}: We use this to detect high-frequency jitter and flickering in rigid monuments like pyramids.
    \item \textbf{LPIPS (Learned Perceptual Image Patch Similarity)}: A deep-learning-based metric that measures perceptual distance. \textit{Justification}: It more accurately reflects human perception of temporal "smoothness" compared to traditional pixel metrics.
    \item \textbf{FVD (Fréchet Video Distance)}: Measures the distributional distance between a set of real videos and our generated samples. \textit{Justification}: This provides a holistic measure of the visual realism and motion quality of the generated Egyptian scenes.
\end{itemize}

\subsection{Ablation Design and Hardware}
We isolate each enhancement in a comprehensive ablation study:
\begin{itemize}
    \item \textbf{Baseline:} Default configuration (Shift 1.0, CFG 4.5, No Neg Prompt).
    \item \textbf{Config A:} Scheduler shift only (1.15).
    \item \textbf{Config B:} Negative prompt only.
    \item \textbf{Config C:} Guidance scale only (6.0).
    \item \textbf{Config D:} All enhancements combined.
\end{itemize}

Each of the 5 configurations is run across all 5 prompts using 3 distinct random seeds (42, 123, 456), yielding 75 total video generations. All experiments generate 60 frames at 30 fps using \texttt{bfloat16} precision. Experiments were executed in parallel on a cluster of NVIDIA RTX A6000 (48GB) GPUs. To maximize inference speed, CPU offloading was strategically managed to retain the Transformer in VRAM during denoising while offloading during the final VAE decode, utilizing \texttt{expandable\_segments} to mitigate memory fragmentation.

\section{Experiments and Results}


\subsection{Quantitative Analysis}
Our results (Table \ref{tab:ablation_results}) demonstrate the efficacy of our zero-cost enhancements. 

The combined configuration (Config D) achieved the highest temporal stability, realizing a significant reduction in LPIPS ($0.0410$ vs $0.0522$ baseline, a 21\% improvement) and the highest average SSIM ($0.9511$). This confirms that the combination of scheduler shifting and negative prompting effectively suppresses the high-frequency jitter and morphing artifacts characteristic of the baseline model.

Regarding semantic alignment, Config C (Guidance Only) achieved the highest overall CLIP-SIM ($33.13$), which is expected given the higher CFG scale. The combined approach (Config D) maintained a highly competitive CLIP-SIM ($32.65$) while drastically outperforming Config C in temporal metrics (LPIPS $0.0410$ vs $0.0575$).

Importantly, our enhancements introduced zero computational overhead. Average inference time for the combined configuration was slightly lower than the baseline (476.3s vs 492.4s), primarily due to optimizations in VRAM allocation.


\begin{figure}[htbp]
\centering
\includegraphics[width=\linewidth]{fig1_metrics_bar.png}
\caption{Comparison of key metrics across different configurations.}
\label{fig:metrics_bar}
\end{figure}

\begin{figure}[htbp]
\centering
\includegraphics[width=\linewidth]{fig2_prompt_heatmap.png}
\caption{Heatmap of semantic alignment and structural preservation across different prompts.}
\label{fig:prompt_heatmap}
\end{figure}

\begin{figure}[htbp]
\centering
\includegraphics[width=\linewidth]{fig3_variance_boxplots.png}
\caption{Variance of metric performance across different random seeds.}
\label{fig:variance_boxplots}
\end{figure}

\subsection{Visual and Distributional Impact}
We analyzed the Fréchet Video Distance (FVD) to measure how our enhancements shifted the generated video distribution relative to the baseline. Config A (Scheduler Only) yielded the lowest FVD ($90.87$), indicating that the scheduler shift preserved the core motion dynamics of the baseline while refining details. Configs involving guidance shifts or negative prompts (B, C, D) resulted in larger FVDs ($~136$-$150$), demonstrating a deliberate divergence from the baseline distribution toward higher-fidelity spaces free of semantic artifacts.

Visual inspection confirms these quantitative findings (see Figure \ref{fig:qualitative_grid}). Baseline videos frequently exhibited "melting" textures on background elements, whereas Config D consistently maintained rigid geometry and sharp feature boundaries across the full 60-frame sequence.

\begin{figure*}[htbp]
\centering
\includegraphics[width=\linewidth,height=0.9\textheight,keepaspectratio]{qualitative_grid.png}
\caption{Qualitative comparison of generated frames showing the baseline versus our combined enhancements (Config D).}
\label{fig:qualitative_grid}
\end{figure*}

\section{Discussion}
The ablation study highlights a critical trade-off in video diffusion models. Increasing CFG (Config C) strongly improves prompt adherence but significantly degrades temporal consistency (worst LPIPS: $0.0575$). However, by coupling high CFG with targeted negative prompting and scheduler optimization (Config D), we successfully mitigate this degradation, achieving best-in-class temporal stability without sacrificing the semantic gains.

\section{Conclusion}
We present a zero-cost, training-free enhancement methodology for open-weight video diffusion models. By applying independent, synergistic modifications to the inference scheduler, text conditioning, and guidance scales, we significantly improve the temporal consistency and structural integrity of generated videos. Our rigorous ablation study confirms that these native edits outperform the baseline model across both perceptual and structural metrics without any additional computational burden. 

\end{document}
